\section{Conclusion}
\label{sec:conclusion}

Polynomial packing makes the witness representation useful beyond
commitment generation: it reduces the relation that the
extension-field prover must check.
For inner and Hadamard products, \methodname encodes a length-\(N\)
relation over \(\Field\) as a polynomial identity and evaluates it
to obtain a length-\(N/\kappa\) relation over \(\ExtensionField\).
The row-polynomial products remain over the base field.
The construction applies to small finite fields of any
characteristic, using distinct field elements as interpolation
nodes.
Under the precomputation and schoolbook assumptions of
Section~\ref{sec:highLevelEfficiency}, with \(N/\kappa\geq\kappa\),
the core reductions require \(O(\kappa N)\) base-field operations
and \(O(N)\) base-field coordinates of prover storage, excluding
commitments and openings.

The packed commitment interface connects these relation reductions
to the original witness. A fixed-basis commitment binds the witness,
while coordinate identities authenticate the evaluations specified
by challenge-dependent packing weights and tensor queries.
Base-field repetition handles the query-side checks, leaving one
terminal opening to an underlying multilinear commitment scheme.
Position-aware permutation arguments then enforce PLONK copy
constraints when packing either a single execution or a batch of
executions of the same circuit.
Theorems~\ref{theorem:single-instance-packing}
and~\ref{theorem:multiple-instance-packing} give the composed
soundness and efficiency bounds, accounting separately for the
intermediate-field labels and underlying commitment costs.

The construction identifies a way to retain small-field arithmetic
without restricting every soundness challenge to the witness field.
For binary arithmetizations with Booleanity constraints, the
three-field instantiation
\(\mathbb F_2\subset\mathbb F_{2^8}\subset\ExtensionField\)
uses byte-field arithmetic for the working-field PIOP steps,
as described in Section~\ref{sec:binary-packing}.
Concrete commitment instantiations and end-to-end evaluation on
quantized inference workloads remain natural next steps.
Extending the current constraint arguments to zero knowledge also
requires compatible masking of the subprotocol transcripts.
These directions build on the same principle: perform the
row-polynomial computation over the small field and reserve the
extension-field proof for the compressed relation.
